% Human source transcription, not native author TeX.
% https://www.jmilne.org/math/CourseNotes/FT500.pdf
% 印刷页45的根式可解定义；页71--72 Proposition 5.26及前置第一方向；页74--75 Lemma 5.32和Theorem 5.33完整双向证明。

\subsection*{Meaning of solvability (source, p. 45)}
For a polynomial $f\in F[X]$, we say that $f(X)=0$ is solvable in radicals if its solutions can be obtained by the algebraic operations of addition, subtraction, multiplication, division, and the extraction of $m$th roots, or, more precisely, if there exists a tower of fields
\[
F=F_0\subset F_1\subset F_2\subset\cdots\subset F_m
\]
such that
\begin{enumerate}[label=(\alph*)]
\item $F_i=F_{i-1}[\alpha_i]$, $\alpha_i^{m_i}\in F_{i-1}$;
\item $F_m$ contains a splitting field for $f$.
\end{enumerate}

\subsection*{Cyclic extensions: the required construction (5.26)}
Let $F$ be a field containing a primitive $n$th root of $1$, some $n\ge2$, and write $\mu_n$ for the group of $n$th roots of $1$ in $F$. Then $\mu_n$ is a cyclic subgroup of $F^\times$ of order $n$ with generator $\zeta$ say.
Consider a field $E=F[\alpha]$ generated by an element $\alpha$ whose $n$th power (but no smaller power) is in $F$. Then $\alpha$ is a root of $X^n-a$, and the remaining roots are the elements $\zeta^i\alpha$, $1\le i\le n-1$. Since these all lie in $E$, $E$ is a Galois extension of $F$, with Galois group $G$ say. For every $\sigma\in G$, $\sigma\alpha$ is also a root of $X^n-a$, and so $\sigma\alpha=\zeta^i\alpha$ for some $i$. Hence $\sigma\alpha/\alpha\in\mu_n$. The map
\[
\sigma\longmapsto\sigma\alpha/\alpha:G\longrightarrow\mu_n
\]
doesn't change when $\alpha$ is replaced by a conjugate, and it follows that the map is a homomorphism:
\[
\frac{\sigma\tau\alpha}{\alpha}
=\frac{\sigma(\tau\alpha)}{\tau\alpha}\frac{\tau\alpha}{\alpha}.
\]
If $\sigma$ lies in the kernel of the map $G\to\mu_n$, then $\sigma\alpha=\alpha$, and so $\sigma$ acts trivially on $E=F[\alpha]$; therefore $\sigma$ is the identity element, and the map is injective. If it is not surjective, then $G$ maps into a subgroup $\mu_d$ of $\mu_n$, some $d\mid n$, $d<n$. In this case, $(\sigma\alpha/\alpha)^d=1$, i.e., $\sigma\alpha^d=\alpha^d$, for all $\sigma\in G$, and so $\alpha^d\in F$, contradicting the hypothesis on $\alpha$. Thus the map is surjective. We have proved the first part of the following statement.
\begin{proposition}[5.26]
Let $F$ be a field containing a primitive $n$th root of $1$. Let $E=F[\alpha]$ where $\alpha^n\in F$ and no smaller power of $\alpha$ is in $F$. Then $E$ is a Galois extension of $F$ with cyclic Galois group of order $n$. Conversely, if $E$ is a cyclic extension of $F$ of degree $n$, then $E=F[\alpha]$ for some $\alpha$ with $\alpha^n\in F$.
\end{proposition}
\begin{proof}
It remains to prove the last statement. Let $\sigma$ generate $G$ and let $\zeta$ generate $\mu_n$. It suffices to find an element $\alpha\in E^\times$ such that $\sigma\alpha=\zeta^{-1}\alpha$, for then $\alpha^n$ is the smallest power of $\alpha$ lying in $F$. As $1,\sigma,\ldots,\sigma^{n-1}$ are distinct homomorphisms $E^\times\to E^\times$, Dedekind's Theorem 5.14 shows that $\sum_{i=0}^{n-1}\zeta^i\sigma^i$ is not the zero function, and so there exists a $\gamma$ such that
\[
\alpha\mathrel{:=}\sum_{i=0}^{n-1}\zeta^i\sigma^i\gamma\ne0.
\]
Now $\sigma\alpha=\zeta^{-1}\alpha$.
\end{proof}

\subsection*{Change of base field (5.32)}
\begin{lemma}[5.32]
Let $f\in F[X]$ be separable, and let $F'$ be a field containing $F$. Then the Galois group of $f$ as an element of $F'[X]$ is a subgroup of the Galois group of $f$ as an element of $F[X]$.
\end{lemma}
\begin{proof}
Let $E'$ be a splitting field for $f$ over $F'$, and let $\alpha_1,\ldots,\alpha_m$ be the roots of $f(X)$ in $E'$. Then $E=F[\alpha_1,\ldots,\alpha_m]$ is a splitting field of $f$ over $F$. Every element of $\operatorname{Gal}(E'/F')$ permutes the $\alpha_i$ and so maps $E$ into itself. The map $\sigma\mapsto\sigma|E$ is an injection
\[
\operatorname{Gal}(E'/F')\longrightarrow\operatorname{Gal}(E/F).
\]
\end{proof}

\subsection*{The counted theorem: Galois's solvability theorem (5.33)}
\begin{theorem}[5.33]
Let $F$ be a field of characteristic $0$. A polynomial in $F[X]$ is solvable if and only if its Galois group is solvable.
\end{theorem}
\begin{proof}
$\Longleftarrow$: Let $f\in F[X]$ have solvable Galois group $G_f$. Let $F'=F[\zeta]$ where $\zeta$ is a primitive $n$th root of $1$ for some large $n$---for example, $n=(\deg f)!$ will do. The lemma shows that the Galois group $G$ of $f$ as an element of $F'[X]$ is a subgroup of $G_f$, and hence is also solvable (GT, 6.6(a)). This means that there is a sequence of subgroups
\[
G=G_0\supset G_1\supset\cdots\supset G_{m-1}\supset G_m=\{1\}
\]
such that each $G_i$ is normal in $G_{i-1}$ and $G_{i-1}/G_i$ is cyclic. Let $E$ be a splitting field of $f(X)$ over $F'$, and let $F_i=E^{G_i}$. We have a sequence of fields
\[
F\subset F[\zeta]=F'=F_0\subset F_1\subset F_2\subset\cdots\subset F_m=E
\]
with $F_i$ cyclic over $F_{i-1}$. Theorem 5.26 shows that $F_i=F_{i-1}[\alpha_i]$ with $\alpha_i^{[F_i:F_{i-1}]}\in F_{i-1}$, each $i$, and this shows that $f$ is solvable.

$\Longrightarrow$: It suffices to show that $G_f$ is a quotient of a solvable group (GT, 6.6(a)). Hence it suffices to find a solvable extension $\widetilde E$ of $F$ such that $f(X)$ splits in $\widetilde E[X]$.
We are given that there exists a tower of fields
\[
F=F_0\subset F_1\subset F_2\subset\cdots\subset F_m
\]
such that
\begin{enumerate}[label=(\alph*)]
\item $F_i=F_{i-1}[\alpha_i]$, $\alpha_i^{r_i}\in F_{i-1}$;
\item $F_m$ contains a splitting field for $f$.
\end{enumerate}
Let $n=r_1\cdots r_m$, and let $\Omega$ be a field Galois over $F$ and containing (a copy of) $F_m$ and a primitive $n$th root $\zeta$ of $1$. For example, choose a primitive element $\gamma$ for $F_m$ over $F$ (see 5.1), and take $\Omega$ to be a splitting field of $g(X)(X^n-1)$ where $g(X)$ is the minimal polynomial of $\gamma$ over $F$. Alternatively, apply 2.9(a).

Let $G$ be the Galois group of $\Omega/F$, and let $\widetilde E$ be the Galois closure of $F_m[\zeta]$ in $\Omega$. According to (3.17(a)), $\widetilde E$ is the composite of the fields $\sigma F_m[\zeta]$, $\sigma\in G$, and so it is generated over $F$ by the elements
\[
\zeta,\alpha_1,\alpha_2,\ldots,\alpha_m,\sigma\alpha_1,\ldots,\sigma\alpha_m,
\sigma'\alpha_1,\ldots.
\]
We adjoin these elements to $F$ one by one to get a sequence of fields
\[
F\subset F[\zeta]\subset F[\zeta,\alpha_1]\subset\cdots
\subset F'\subset F''\subset\cdots\subset\widetilde E
\]
in which each field $F''$ is obtained from its predecessor $F'$ by adjoining an $r$th root of an element of $F'$ ($r=r_1,\ldots,r_m$, or $n$). According to (5.8) and (5.26), each of these extensions is abelian (and even cyclic after the first), and so $\widetilde E/F$ is a solvable extension.
\end{proof}
